\documentclass[12pt]{article}

\usepackage{bbm}
\usepackage{fancyhdr}
\usepackage{titling}
\usepackage{enumitem}
\usepackage{amsmath}
\usepackage{dsfont}
\usepackage{amsfonts}
\usepackage{amssymb}
\usepackage{amsthm}
\usepackage{geometry}
\usepackage{tikz}
\usepackage{algorithm}
\usepackage{algpseudocode}
\geometry{verbose,tmargin=1in,bmargin=1in,lmargin=1in,rmargin=1in}
\renewcommand{\vec}[1]{\mathbf{#1}}
\newcommand{\dec}{\text{dec}}
\newcommand{\divi}{\mathbin{\textnormal{div}}}
\title{Homework 1}
\author{Satvik Prasad}
\date{}
\newtheorem{lemma}{Lemma}[section]
\fancypagestyle{plain}{%
  \fancyhf{}
  \fancyhead[L]{\theauthor}
  \fancyhead[R]{\thedate}
}
\pagestyle{plain}
\makeatletter
\def\@maketitle{%
  \newpage
  \null
  \vskip 1em%
  \begin{center}%
    \let \footnote \thanks
    {\LARGE \@title \par}%
    \vskip 1em%
  \end{center}%
  \par
\vskip 1em}
\makeatother
\usepackage{cmbright}

\newcommand{\mat}[1]{\mathbf{#1}}

% ---------- Probability theory ----------
% Core operators (optional argument-free; use as \E{X}, \Var{X}, \Cov{X}{Y})
\DeclareMathOperator{\EOp}{\mathbb{E}}
\DeclareMathOperator{\VarOp}{Var}
\DeclareMathOperator{\CovOp}{Cov}
\DeclareMathOperator{\CorrOp}{Corr}

\renewcommand{\P}[1]{\mathbb{P}\left(#1\right)}        % \P{A}
\newcommand{\E}[1]{\EOp\left[#1\right]}                % \E{X}
\newcommand{\Var}[1]{\VarOp\left(#1\right)}            % \Var{X}
\newcommand{\Cov}[2]{\CovOp\left(#1,\,#2\right)}       % \Cov{X}{Y}
\newcommand{\Corr}[2]{\CorrOp\left(#1,\,#2\right)}     % \Corr{X}{Y}
\newcommand{\SD}[1]{\operatorname{SD}\left(#1\right)}  % \SD{X}

% Subscripted versions, e.g. \Esub{\theta}{X}, \Psub{\theta}{A}
\newcommand{\Esub}[2]{\EOp_{#1}\left[#2\right]}
\newcommand{\Psub}[2]{\mathbb{P}_{#1}\left(#2\right)}

% Conditioning, indicators, independence
\newcommand{\given}{\,\middle|\,}                      % \E{X \given Y}
\newcommand{\ind}[1]{\mathbbm{1}\left\{#1\right\}}     % \ind{X > 0}
\newcommand{\indep}{\perp\!\!\!\perp}                  % X \indep Y
\newcommand{\iid}{\overset{\text{iid}}{\sim}}          % X_i \iid \Normal(0,1)

% Convergence
\newcommand{\convp}{\xrightarrow{\;p\;}}               % in probability
\newcommand{\convd}{\xrightarrow{\;d\;}}               % in distribution
\newcommand{\convas}{\xrightarrow{\text{a.s.}}}        % almost surely
\newcommand{\convLp}[1]{\xrightarrow{L^{#1}}}          % in L^p

% Probability spaces
\newcommand{\F}{\mathcal{F}}                           % sigma-algebra
\newcommand{\B}{\mathcal{B}}                           % Borel sigma-algebra
\newcommand{\pspace}{(\Omega, \F, \mathbb{P})}

% Common distributions
\newcommand{\Normal}{\mathcal{N}}
\newcommand{\Bern}{\operatorname{Bernoulli}}
\newcommand{\Bin}{\operatorname{Binomial}}
\newcommand{\Pois}{\operatorname{Poisson}}
\newcommand{\Unif}{\operatorname{Uniform}}
\newcommand{\Exp}{\operatorname{Exp}}
\newcommand{\Geom}{\operatorname{Geometric}}
\newcommand{\Gam}{\operatorname{Gamma}}
\newcommand{\Beta}{\operatorname{Beta}}

% Number sets
\newcommand{\R}{\mathbb{R}}
\newcommand{\N}{\mathbb{N}}
\newcommand{\Z}{\mathbb{Z}}


\begin{document}
\begin{enumerate}
  \item
    \begin{enumerate}
      \item Consider the case where we have a collection of coins $C
        = [4, 3, 1],$ and we have $V = 14.$ The greedy approach will
        select $4 \times 3$ and $1 \times 2$, i.e, $5$ coins, whereas
        we can choose $4 \times 2$, $3 \times 2$ to construct $V =
        14$ with 4 coins, meaning the greedy approach is suboptimal.

      \item
        \begin{enumerate}
          \item Construct $T[i]$ where $T[i]$ is the minimum number
            of coins required to attain $V = i.$
          \item We have
            $$
            T[j] = \min_{c \in C, \ c \leq j}(T[j - c]) + 1.
            $$

            \textbf{Proof}: the true minimal way to construct $V = j$
            will contain some final chosen coin, $c.$ If the
            remaining coins are not minimal for $V = j - c,$ then we
            can construct a more minimal selection for $V = j.$
            Hence, given we know $T[j -c ]$ is correct for all
            possible coins $c$, $T[j]$ will be correct.

          \item The base case is simply $T[0] = 0,$ since there are
            $0$ coins required to form $V = 0$ from the set of coins
            (assuming coin values are all positive).

          \item Cell $T[V]$ contains the solution.

          \item We must perform the recursion $V$ times, and each
            recursion has $\mathcal O(|C|)$ iterations, meaning our
            algorithm runs in $\mathcal O(V|C|)$ time.
        \end{enumerate}
    \end{enumerate}

  \item
    \begin{enumerate}
      \item $A = [3, 1, 1, 1, 1, 1, 1]$ with $k = 1$ can be entirely
        protected if we place a guard at entry $i = 4,$ but the
        greedy algorithm will choose at least $i = 1,$ and yield a
        score of only $6$ (as opposed to $9$).

      \item
        \begin{enumerate}
          \item Define $T[i, j]$ as the maximum score achievable with
            at most $i$ knights when the gallery consists only of
            paintings $1, \ldots, j.$

          \item For $i \geq 1,$ our recurrence is given by
            \begin{align*}
              T[i, j] =
              \begin{cases}
                \max\left\{
                  T[i, j - 1], \
                  C[j - 3] + T[i - 1, j - 7]
                \right\} & j \geq 7 \\[2mm]
                \sum_{t = 1}^{j} A[t]  & 1 \leq j \leq 6
              \end{cases}
            \end{align*}
            where $C[p] = \sum_{t = p - 3}^{p + 3}A[t]$ represents
            the sum of all paintings protected by a knight placed at
            index $p.$ Note that $C[p]$ is only ever used with $4
            \leq p \leq n - 3,$ so the sum never leaves the array.

            \paragraph{Proof} If $1 \leq j \leq 6,$ a single knight
            at position $\max(1, j - 3)$ protects every painting in
            $1, \ldots, j,$ so the score is the sum of the whole prefix.

            Now let $j \geq 7,$ and consider an optimal placement of
            at most $i$ knights over paintings $1, \ldots, j.$ Either
            painting $j$ is not protected, or it is.

            In the first case, no knight stands on positions $j - 3,
            \ldots, j,$ so the same placement is valid for paintings
            $1, \ldots, j - 1$ with the same score, which is at most
            $T[i, j - 1].$

            In the second case, some knight $g$ stands on a position
            in $j - 3, \ldots, j.$ Moving $g$ to position $j - 3$
            protects paintings $j - 6, \ldots, j,$ which is a
            superset of what it protected before, so the score does
            not decrease. Any other knight that protects a painting
            in $j - 6, \ldots, j$ can be moved left to position at
            most $j - 7$ without losing any painting with index $\leq
            j - 7$ that it protected. The remaining knights then form
            a placement of at most $i - 1$ knights over paintings $1,
            \ldots, j - 7.$ If this placement were not maximal,
            replacing it with a maximal one and keeping $g$ at
            position $j - 3$ would give a better placement, a
            contradiction. Hence the score is $C[j - 3] + T[i - 1, j - 7].$

            Thus, if $T[i - 1, j - 7]$ and $T[i, j - 1]$ are correct,
            then $T[i, j]$ is correct.

          \item Our base cases are $T[0, j] = 0$ for $0 \leq j \leq
            n,$ and $T[i, 0] = 0$ for $0 \leq i \leq k.$

          \item Cell $T[k, n]$ contains the solution, since this is
            the maximum score with at most $k$ knights over all
            paintings $1, \ldots, n.$

          \item There are $\mathcal O(nk)$ cells and each cell
            requires $\mathcal O(1)$ work if we pre-compute $C$ and
            the prefix sums of $A,$ meaning this algorithm takes
            $\mathcal O(nk)$ time.
        \end{enumerate}
    \end{enumerate}

  \item
    \begin{enumerate}
      \item Define $T[i, j, 0]$ as the maximum achievable rating with
        at most $i$ dollars, where the largest episode is $j,$ and $j
        - 1$ is also chosen. Define $T[i, j, 1]$ as the equivalent,
        but where $j$ is chosen and $j - 1$ is not chosen. Define
        $T[i, j, 2]$ as the equivalent, but the largest episode has index $< j.$

      \item Our recursion is defined as
        \begin{align*}
          T[i, j, 0] &= \max(T[i - 1, j - 1, 0], T[i, j - 1, 1]) + R[j], \\
          T[i, j, 1] &= T[i - 1, j - 1, 2] + R[j], \\
          T[i, j, 2] &= \max(T[i, j - 1, 0], T[i, j - 1, 1], T[i, j -1, 2])
        \end{align*}
        if $i \geq 1, j \geq 2.$

        Any optimal movie selection has some largest movie $j.$ Then,
        if $j - 1$ is also in this selection, and $j - 2$ is not,
        adding $j$ is free. Otherwise (whether $j - 1$ is not in the
        selection, or $j - 1$ is and $j - 2$ is not), adding $j$ adds
        $1$ to the cost of the set of movies.

      \item Our base cases are as follows: whenever $i = 0$ or $j = 0$,
        \begin{align*}
          T[i, j, 2] &= 0, \\
          T[i, j, 0] = T[i, j, 1] &= -\infty.
        \end{align*}
        The recursion then applies for all $i \geq 1, \ j \geq 1.$

      \item We have $\mathcal O(nk)$ cells that each require
        $\mathcal O(1)$ time to compute, meaning our algorithm runs
        in $\mathcal O(nk)$ time.
    \end{enumerate}

  \item
    \begin{enumerate}
      \item
        \begin{enumerate}
          \item We define $T[i]$ as the minimum total badness
            achievable by inserting line breaks for the substring
            consisting of words from $1 \ldots i.$

          \item Define the recursion as
            $$
            T[i] = \min_{1 \leq j \leq i}(\texttt{Badness}(j, i) + T[j - 1]),
            $$
            when $i \geq 1.$

            This is correct because every single partitioning can be
            decomposed into the final partition $p$ and all previous
            partitions $S$. If $S$ not minimal, then taking the
            minimal partitioning instead of $S$ will produce a more
            minimal partitioning. Hence, we only need to iterate over
            the final partition and the optimal partition preceding
            (which exists in our table for smaller $i$).

          \item Base case consists of $T[0] = 0.$

          \item $\mathcal O(n)$ iterations, where each cell requires
            $\mathcal O(n)$ time to compute, yielding $\mathcal O(n^2)$ time.
        \end{enumerate}

      \item
        \begin{algorithm}
          \begin{algorithmic}
            \Procedure{ExtractIdx}{$T, n$}
            \State $B \gets \texttt{vector}()$

            \State $j \gets n$

            \vspace{3mm}

            \While{$j \geq 1$}
            \State $i \gets j$

            \vspace{3mm}

            \While{$T[j] - T[i - 1] \neq \texttt{Badness}(i, j)$}
            \State $i \gets i - 1$
            \EndWhile

            \vspace{3mm}

            \If{$i = 1$}
            \State \textbf{break}
            \EndIf

            \vspace{3mm}

            \State $B.\texttt{push\_back}(i - 1)$ \Comment{Final line
            starts at word $i$}

            \State $j \gets i - 1$
            \EndWhile

            \vspace{3mm}

            \State \Return $\texttt{reversed}(B)$
            \EndProcedure
          \end{algorithmic}
        \end{algorithm}
    \end{enumerate}
\end{enumerate}
\end{document}
